\documentclass[11pt,aspectratio=169]{beamer}
\usepackage{hp-dfr-slides}
\graphicspath{{../../../assets/}}

\title{Goal-oriented DFR}
\subtitle{A primal and an adjoint network, trained together}
\author{Carlos Vera-Ciro}
\institute{Goal-Oriented DFR series \textbar\ Part 2}
\date{}

\begin{document}

\begin{frame}[plain,noframenumbering]
  \titlepage
\end{frame}

% =============================================================================
\begin{frame}{Where we are}
  \begin{itemize}\itemsep0.6em
    \item Part 1: DFR gives a loss, the $H^{-1}$ dual norm of the residual, that
          is equivalent to the global energy error.
    \item But we often want a single \important{quantity of interest} $J(u)$, and
          global control is indirect and wasteful.
    \item This deck: weight the DFR loss toward $J$, using the dual-weighted
          residual (DWR) idea.
  \end{itemize}
\end{frame}

% =============================================================================
\section{The dual-weighted residual idea}
% =============================================================================

\begin{frame}{The adjoint as an importance map}
  Classical finite elements answer ``where does residual matter for my QoI?''
  with an \important{adjoint} solution $z^*$.
  \medskip
  \begin{itemize}\itemsep0.5em
    \item $z^*$ is large where a residual error most affects $J(u)$, and small
          where it does not.
    \item The QoI error is then the residual \emph{weighted by that map}:
          \[
            J(u^*) - J(u_h) \;=\; \Res(u_h)(z^*) .
          \]
  \end{itemize}
  \takeaway{Weight the residual by the adjoint, and you localize the error onto
  the quantity you care about. Part 3 proves this identity.}
\end{frame}

% =============================================================================
\section{The construction}
% =============================================================================

\begin{frame}{Two networks}
  Train a \important{primal} network and an \important{adjoint} network together.
  \medskip
  \begin{center}
  \begin{tikzpicture}[node distance=7mm and 14mm]
    \node[accentnode] (u) {primal $u_\theta$};
    \node[accentnode, below=of u] (z) {adjoint $z_\phi$};
    \coordinate (mid) at ($(u)!0.5!(z)$);
    \node[opnode, right=16mm of mid] (pair) {$\bigl|\inner{\Res(u_\theta)}{z_\phi}\bigr|$\\QoI-weighted loss};
    \draw[localarrow] (u.east) -- (pair.west);
    \draw[localarrow] (z.east) -- (pair.west);
  \end{tikzpicture}
  \end{center}
  \medskip
  \begin{itemize}\itemsep0.4em
    \item The primal minimizes the QoI-weighted residual, putting its accuracy
          where the adjoint says it matters.
    \item The adjoint solves its \emph{own} DFR problem, so it genuinely
          approximates $z^*$.
  \end{itemize}
\end{frame}

% -----------------------------------------------------------------------------
\begin{frame}{The subtlety that makes it work: a stop-gradient}
  The goal term $\bigl|\inner{\Res(u_\theta)}{z_\phi}\bigr|$ is a single scalar.
  If the \important{adjoint} were allowed to descend it too:
  \medskip
  \begin{itemize}\itemsep0.5em
    \item[\alert{$\times$}] the optimizer would drive the pairing to zero
          trivially, by making $z_\phi$ \emph{orthogonal} to the residual,
    \item[\alert{$\times$}] without $z_\phi$ approximating the true adjoint at
          all. The localization is lost.
  \end{itemize}
  \medskip
  \good{Fix:} hold the adjoint fixed in the goal term (a stop-gradient). The
  adjoint is trained \emph{only} by its own residual $\norm{\Resad(z_\phi)}_{U^*}^2$;
  the primal is guided by the QoI-weighted term.
  \takeaway{Stop-gradient stops the adjoint from cheating. The difference between
  a working method and a degenerate one.}
\end{frame}

% -----------------------------------------------------------------------------
\begin{frame}{The training objective}
  \[
    \mathcal{L}(\theta,\phi) =
      \underbrace{\bigl|\Res(u_\theta)(\bar{z}_\phi)\bigr|}_{\text{goal term (primal only)}}
      \;+\; \beta\,\underbrace{\norm{\Res(u_\theta)}_{V^*}^2}_{\text{primal DFR}}
      \;+\; \alpha\,\underbrace{\norm{\Resad(z_\phi)}_{U^*}^2}_{\text{adjoint DFR}}
  \]
  \medskip
  \begin{itemize}\itemsep0.5em
    \item $\bar{z}_\phi$: adjoint held fixed (the stop-gradient).
    \item The pairing $\Res(u_\theta)(z_\phi) = \int_\Omega (f + \Delta u_\theta)\,
          z_\phi\, dx$ is evaluated on the same grid DFR already uses; no extra
          quadrature.
    \item Same spectral evaluation and boundary-cutoff construction as plain DFR.
          The only new ingredients are the adjoint and the goal term.
  \end{itemize}
\end{frame}

% -----------------------------------------------------------------------------
\begin{frame}{The algorithm, in one loop}
  \begin{enumerate}\itemsep0.5em
    \item Precompute the DST matrix and the $H^{-1}$ weights.
    \item Each step, by DST, form the residual coefficients $\hat{\Res}(u_\theta)$
          and $\hat{\Resad}(z_\phi)$.
    \item Primal loss $\mathcal{L}_u = |\Res(u_\theta)(\bar{z}_\phi)|
          + \beta\,\norm{\Res(u_\theta)}_{V^*}^2$ updates $\theta$.
    \item Adjoint loss $\mathcal{L}_z = \alpha\,\norm{\Resad(z_\phi)}_{U^*}^2$
          updates $\phi$.
    \item Adam, then an L-BFGS polish.
  \end{enumerate}
  \medskip
  Boundary conditions are enforced by a cutoff: $u_\theta = \varphi\,\tilde{u}_\theta$,
  with $\varphi$ vanishing on $\partial\Omega$; likewise for $z_\phi$.
  \takeaway{Next: what this construction lets us prove.}
\end{frame}

\end{document}
